\subsection{A Primitive Condition for a Risk-Specific Adjustment Option}
\label{sec:r7_primitive}
The option decomposition identifies the quantity that changes a risk ranking; a primitive condition is needed to predict its sign. Consider a two-stage specialization that isolates the preference channel. A risk class $\sigma\in\{+,-\}$ delivers a positive consumption lottery $C_\sigma$ and an exogenously specified discounted operating duration $A_\sigma$. Before consumption is realized, the agent chooses $\theta\in[-\bar\theta,\bar\theta]$ at cost $k\theta^2/2$. Set
\[
 F_\sigma(\theta)=\E\frac{C_\sigma^{1-u_0-\theta}}{1-u_0-\theta},\quad
 B_\sigma=F_\sigma(0),\quad g_\sigma=F'_\sigma(0),\quad
 M_\sigma=\max_{|\theta|\leq\bar\theta}[-F''_\sigma(\theta)],
\]
where $u_0>1+\bar\theta$ and consumption has finite positive support. The class value is $B_\sigma+dA_\sigma+\Omega_\sigma(k)$, with
$\Omega_\sigma=\max_{|\theta|\leq\bar\theta}\{F_\sigma(\theta)-B_\sigma-k\theta^2/2\}$.
Unlike the optimized-policy moments in \eqref{eq:risk_derivatives}, $g_\sigma$ and $M_\sigma$ are computed directly from the consumption lottery and preference index.

\begin{proposition}[Consumption exposure and the adjustment premium]
\label{prop:r7_primitive}
For $u>1$,
\begin{equation}
 U_{uu}(c,u)=-\frac{c^{1-u}}{(u-1)^3}
 \left\{[1+(u-1)\log c]^2+1\right\}<0.
 \label{eq:r7_curvature}
\end{equation}
If $|g_\sigma|/(k+M_\sigma)\leq\bar\theta$, then
\begin{equation}
 \frac{g_\sigma^2}{2(k+M_\sigma)}\leq\Omega_\sigma(k)
 \leq\frac{g_\sigma^2}{2k}.
 \label{eq:r7_option_bounds}
\end{equation}
Consequently, if $A_+-A_-=:D_A>0$ and
\begin{equation}
 \underline\Omega:=\frac{g_+^2}{2(k+M_+)}-\frac{g_-^2}{2k}>0,
 \label{eq:r7_primitive_condition}
\end{equation}
preference adjustment lowers the risk-indifference contract by at least
$\underline\Omega/D_A$. Specifically, with $d^0=(B_--B_+)/D_A$, every
$d\in(d^0-\underline\Omega/D_A,d^0)$ strictly favors the nonpositive class without adjustment and the positive class with adjustment.
\end{proposition}

The argument uses global concavity in the preference index, not an interior first-order condition for an optimized portfolio. Low consumption affects the adjustment incentive through
$U_u(c,u)=c^{1-u}[1+(u-1)\log c]/(u-1)^2$.
A downside consumption exposure can therefore increase the magnitude of the preference-adjustment benefit even when its unadjusted utility is lower. Condition \eqref{eq:r7_primitive_condition} compares that gain with curvature and effort cost; greater risk by itself is not sufficient. Supplement S.11 proves the bounds and shows how $M_\sigma$ can be bounded using endpoint preferences rather than an unverified curvature grid.

For a declared family, take $u_0=2$, $\bar\theta=0.2$, $C_-=0.5$, and
$C_+=0.05$ with probability $p$ and $0.8$ otherwise, with $D_A=0.5$.
For every $p\in[0.06,0.10]$ and $k\in[20,80]$, the primitive bound gives
$\Omega_+-\Omega_->0.00186$ and a threshold reduction exceeding $0.00372$ in operating-benefit units. This is a continuum prediction from the primitive inequalities, not interpolation of a nine-point exercise. The finite support and preference interval imply $|g_+|\geq1.4820$, $|g_-|\leq0.6138$, and $M_+\leq17.342$; these conservative bounds imply the stated uniform result. The larger, point-specific lower bounds and the exactly optimized options are reported in the supplement. This specialization explains a channel; it is not offered as a calibration or approximation of the full stopped diffusion. The full finite economy is examined separately with decision-specific certificates.

\subsection{Certifying a Risk Ranking over a Joint Counterfactual Region}
\label{sec:r7_decision}
A welfare tolerance is not a tolerance for the sign of a small class-value difference. We therefore construct upper and lower values separately for each risk class and each adjustment regime. Fix $I_\lambda=[a,b]$ and $I_d=[d_l,d_h]$ and write $t=(\lambda-a)/(b-a)$ and $s=(d-d_l)/(d_h-d_l)$. For a fixed policy, its value is polynomial in $t$ and affine in $d$. For a fixed $\lambda$, the optimal class value is convex in $d$, since the feasible policy set is unchanged and every policy value is affine.

Let $U_\sigma^l(t)$ and $U_\sigma^h(t)$ be valid transition-polynomial upper values at the two contracts, constructed using the constrained count or corrected-chord recursion. Then
\[
 U_\sigma(t,s)=(1-s)U_\sigma^l(t)+sU_\sigma^h(t)
\]
is an upper value throughout the rectangle. Evaluate any feasible class-policy bank independently to obtain lower polynomials $L_{\sigma,p}(t,s)$. Denote their tensor Bernstein coefficients on a cell by $U_{\sigma,ij}$ and $L_{\sigma,p,ij}$.

\begin{proposition}[Decision-specific tensor certificate]
\label{prop:r7_tensor}
Suppose the upper and lower polynomial constructions, including their coefficient operations, each have absolute error at most $e$. Throughout the closed cell,
\begin{align}
 \underline\Delta&:=\max_{p\in\mathcal B_+}\min_{i,j}
   \{L_{+,p,ij}-U_{-,ij}\}-2e\leq\Delta,\label{eq:r7_tensor_sign}\\
 \Delta&\leq\overline\Delta:=\min_{p\in\mathcal B_-}\max_{i,j}
   \{U_{+,ij}-L_{-,p,ij}\}+2e.\nonumber
\end{align}
In particular, $\underline\Delta^{\rm adj}>0$ and $\overline\Delta^0<0$ certify opposite risk rankings under the same contract and shock-law parameter. They also imply
$\Omega_+-\Omega_-\geq\underline\Delta^{\rm adj}-\overline\Delta^0>0$.
\end{proposition}
\begin{proof}
A tensor Bernstein polynomial is a convex combination of its coefficients. Each positive-class policy is a feasible lower value and each negative-class upper polynomial bounds that class optimum; their difference therefore bounds $\Delta$ from below. Taking a coefficient minimum, followed by a maximum over feasible policies, preserves that bound. Interchanging the classes gives the upper inequality. The two absolute error allowances add. Applying the option decomposition gives the final statement.
\end{proof}

This certificate does not require uniform duration separation, a regular active-policy cell, or uniqueness of a threshold. Its object is a strict choice on a declared region. It uses the same transition parameter in all periods, permits later risky positions to change sign, and imposes the adjustment restriction at every date. In the application below, the class-value tolerances are generated by these coefficient comparisons rather than borrowed from the $10^{-3}$ generic welfare result.
